% Nicolò Ciccarone — Curriculum (italiano)
% Stesso impianto del PDF inglese: Cabin, Cochineal, Inconsolata, Font Awesome 5.
\documentclass[10pt,a4paper]{article}

\usepackage[margin=0.9cm]{geometry}
\usepackage{iftex}
\ifPDFTeX
  \usepackage[T1]{fontenc}
  \usepackage[utf8]{inputenc}
\fi
\usepackage[italian]{babel}
\usepackage[scale=.95]{cabin}
\usepackage{cochineal}
\ifPDFTeX\else
\setmonofont{Inconsolatazi4}[
  Scale=0.9,
  Extension=.otf,
  UprightFont=*-Regular,
  BoldFont=*-Bold]
\fi
\usepackage{fontawesome5}
\usepackage{xcolor}
\usepackage{enumitem}
\usepackage{tikz}
\usepackage{hyperref}

\definecolor{accent}{rgb}{0.57098,0.02196,0.19766}
\definecolor{rulecolor}{rgb}{0.72,0.8047,0.42589}

\hypersetup{
  hidelinks,
  pdftitle={Nicolò Ciccarone},
  pdfauthor={Nicolò Ciccarone}
}
\urlstyle{tt}

\pagestyle{empty}
\setlength{\parindent}{0pt}
\setlength{\parskip}{0pt}
\setlength{\emergencystretch}{1em}
\setlist[itemize]{
  leftmargin=1em,
  itemsep=2pt,
  topsep=3pt,
  parsep=0pt,
  partopsep=0pt,
  label=\textbullet,
  before=\footnotesize
}

\newcommand{\cvsection}[1]{%
  \par\vspace{9pt}%
  \noindent\begin{tikzpicture}[line width=0pt]
    \shade[left color=rulecolor, right color=white]
      (0,0) rectangle (\textwidth,2.5pt);
    \node[anchor=base west, inner xsep=1.2pt, inner ysep=0pt, outer sep=0pt,
          font=\sffamily\bfseries\Large] at (0,4.33pt) {#1};
  \end{tikzpicture}\par\vspace{5pt}%
}
\newcommand{\cvsub}[1]{%
  \par\vspace{2pt}%
  {\sffamily\bfseries\large #1\par}\vspace{3pt}%
}

% Layout a colonne condiviso da tutte le sezioni: data | icona | contenuto
\newlength{\datecol}
\newlength{\colgap}
\newlength{\iconw}
\newlength{\contentw}
\settowidth{\datecol}{\small Novembre 2022 -- dicembre 2023}
\setlength{\colgap}{0.45cm}
\setlength{\iconw}{16pt}
\setlength{\contentw}{\dimexpr\textwidth-\datecol-\colgap-\iconw\relax}

\newcommand{\cvicon}{\textcolor{accent}{\small\faBookmark}}
\newcommand{\cvitem}[3][6pt]{%
  \par\noindent
  \makebox[\datecol][r]{\small #2}\hspace{\colgap}%
  \makebox[\iconw][l]{\cvicon}%
  \parbox[t]{\contentw}{#3}\par
  \vspace{#1}%
}
\newcommand{\edudesc}[1]{\\{\footnotesize #1\par}}

\newcommand{\role}[1]{\textbf{#1}}
\newcommand{\org}[1]{\textbf{\textit{#1}}}

\begin{document}
{\sffamily\bfseries\LARGE Nicolò Ciccarone\par}
\vspace{1.5pt}
{\small
\textcolor{accent}{\faPhone}\hspace{0.45em}+39 3663221036%
\hspace{1.55em}%
\textcolor{accent}{\faEnvelope[regular]}\hspace{0.45em}\texttt{niciccarone@gmail.com}\par
\vspace{1.2pt}
\textcolor{accent}{\faLinkedin}\hspace{0.4em}\href{https://www.linkedin.com/in/nicolociccarone1997/}{\texttt{https://www.linkedin.com/in/nicolociccarone1997/}}%
\hspace{1.35em}%
\textcolor{accent}{\faGithub}\hspace{0.4em}\href{https://github.com/nclcc}{\texttt{https://github.com/nclcc}}\par
}

\cvsection{Esperienza professionale}

\cvitem[8pt]{Gennaio 2025 -- in corso}{%
  \role{Data Scientist / Data Engineer, \org{Complexa Labs}, Bari (IT)}
  \begin{itemize}
    \item Realizzazione di soluzioni di analytics come consulente freelance per clienti enterprise (fatturato superiore a 1 milione di euro).
    \item Sviluppo di un motore ATS di time-to-fill: modello di sopravvivenza a tempo discreto score-driven (GAS) su un pannello giornaliero a livello di posizione vacante, con posizioni aperte censurate a destra e covariate giornaliere del funnel di selezione, in Polars, messo in produzione su SageMaker.
    \item Reingegnerizzazione delle pipeline in operazioni Spark-native (PySpark): join nativi, funzioni finestra e broadcast hints, per ottimizzare le prestazioni.
    \item Sviluppo di un simulatore di scenari macroeconomici, con orchestrazione distribuita delle pipeline, per prevedere gli impatti sul business al variare delle condizioni di mercato.
    \item Sviluppo di un modello di Customer Lifetime Value (CLTV) che combina analisi di sopravvivenza mixture-cure, regressione del margine con gradient boosting e segmentazione calibrata, per la previsione dei ricavi e il targeting del Next Best Product.
    \item Sviluppo di algoritmi NLP per l'automazione dell'elaborazione documentale, messi in produzione come microservizi FastAPI in Python.
  \end{itemize}}
\cvitem[8pt]{Gennaio 2024 -- novembre 2024}{%
  \role{Data Scientist, \org{acmetric B.V.}, Amsterdam (NL)}
  \begin{itemize}
    \item Sviluppo e manutenzione di un motore di pricing e previsione della domanda per Adidas International B.V., basato su modelli di Causal AI.
    \item Progettazione e messa in produzione di un modello di previsione di churn e contenzioso per LATAM Airlines, con Gradient Boosting (GBM).
  \end{itemize}}
\cvitem[0pt]{Novembre 2022 -- dicembre 2023}{%
  \role{Data Scientist, \org{MOOS Solutions B.V.}, Amsterdam (NL)}
  \begin{itemize}
    \item Sviluppo e manutenzione di un motore di intelligenza artificiale per gli insight in tempo reale sull'inventario, basato su tecniche di Mixture of Experts (MoE).
  \end{itemize}}

\cvsection{Formazione}

\cvitem{Luglio 2023}{%
  \textbf{Summer School in Deep Learning, Tinbergen Institute, Amsterdam}
  \edudesc{Scuola estiva intensiva di livello dottorale, tenuta da Eran Raviv, PhD, sul deep learning, con un focus sull'\textbf{intelligenza artificiale generativa}.}}
\cvitem{2021 -- 2022}{%
  \textbf{M.Sc.\ in Econometria e Ricerca Operativa, Vrije Universiteit Amsterdam}
  \edudesc{Tesi: \textit{Volatilità dei prezzi ed estremi climatici: distinguere il segnale dal rumore} (\textit{Price volatility and climatic extremes: distinguishing signal from noise}). Discussa il 18 luglio 2022. \textit{Voto finale}: 8,0/10.}}
\cvitem{2019 -- 2021}{%
  \textbf{Laurea magistrale in Economia, Università di Bologna, curriculum Econometria}
  \edudesc{Tesi: \textit{Copertura contro il cambiamento climatico: un approccio con processi di Hawkes} (\textit{Hedging against climate change: A Hawkes process approach}). Discussa il 23 luglio 2021. \textit{Voto finale}: 110/110 e lode.}}
\cvitem[0pt]{2015 -- 2018}{%
  \textbf{Laurea triennale in Economia e Management, Università di Trento}
  \edudesc{Ammissione all'\textit{Honours Programme} del Dipartimento di Economia (riservato agli studenti particolarmente meritevoli). \textit{Voto finale}: 110/110 e lode.}}

\cvsection{Competenze}

\cvitem[3pt]{Lingue}{\raggedright
  \role{Italiano} (\textit{madrelingua}), \role{Inglese} (\textit{C1/C2}), \role{Francese} (\textit{B1}), \role{Tedesco} (\textit{A2})}
\cvitem[3pt]{IT e programmazione}{\raggedright
  \role{Python} (\textit{esperto}), \role{R} (\textit{esperto}), \textbf{\LaTeX{}} (\textit{esperto}), \role{SQL} (\textit{avanzato}), \role{AWS} (\textit{avanzato}), \role{Docker} (\textit{avanzato}), \role{Microsoft Office} (\textit{avanzato}), \role{strumenti GIS} (\textit{avanzato}), \role{Kubernetes} (\textit{intermedio}), \role{Spark} (\textit{intermedio}), \role{Databricks} (\textit{intermedio}), \role{PyTorch} (\textit{intermedio}), \role{FastAPI} (\textit{intermedio}), \role{Celery} (\textit{intermedio}), \role{MongoDB} (\textit{intermedio}), \role{GCP} (\textit{intermedio}), \role{Microsoft Azure} (\textit{intermedio}), \role{Grafana} (\textit{intermedio}), \role{Julia} (\textit{base})}
\cvitem[0pt]{Servizi cloud}{\raggedright
  \role{AWS}: SageMaker, ECS, EKS, ECR, CloudWatch\enspace---\enspace\role{Azure}: Pipelines, ML, AKS, ACI, ACR, Monitor, Fabric\enspace---\enspace\role{GCP}: Vertex AI Workbench, Colab Enterprise}

\cvsection{Esperienze rilevanti}

\cvsub{Certificazioni e corsi}
\cvitem[3pt]{2021}{\role{TOEFL}. Punteggio complessivo: 107/120, livello C1 (QCER).}
\cvitem[0pt]{2016}{\role{Livello A1, corso estivo intensivo di tedesco}, presso lo Sprachzentrum di Göttingen (DE).}

\cvsub{Riconoscimenti}
\cvitem[3pt]{2019}{\role{INPS (Istituto Nazionale della Previdenza Sociale)}, borsa di studio per merito, edizione 2019.}
\cvitem[0pt]{2018}{\role{Premi dell'Università di Trento}, borsa di studio per merito, edizione 2018.}

\end{document}
